% Канонический TikZ-рисунок варианта 16. Править только здесь.
\newcommand{\figCyclicOpposite}[1]{%
\begin{tikzpicture}[scale=0.78,line cap=round,line join=round,every node/.style={font=\small}]
  \def\R{2.3}
  \coordinate (A) at (190:\R); \coordinate (B) at (130:\R);
  \coordinate (C) at (0:\R); \coordinate (D) at (236:\R);
  \draw[line width=0.9pt] (0,0) circle (\R);
  \draw[line width=0.9pt] (A)--(B)--(C)--(D)--cycle;
  \pic[draw,line width=0.65pt,angle radius=5mm] {angle=D--A--B};
  \node at (-1.37,-.10) {\(#1^\circ\)};
  \node[left] at (A) {\(A\)}; \node[above left] at (B) {\(B\)};
  \node[right] at (C) {\(C\)}; \node[below left] at (D) {\(D\)};
\end{tikzpicture}}
